\section{Claim-by-Claim Integrated Argument}
\label{sec:oc133-claim-by-claim-integrated-argument}
The following integrated argument turns the claim and theorem registers into a reader-facing scientific route. It is intentionally longer than a manifest because a publication-grade monograph must explain why each promoted claim exists, what problem it solves, which evidence carries it, and where the claim stops. The register remains available for audit; the argument below is written for scientific reading.

\subsection{Shared Claim Reading Method}
Every promoted model-core claim is read by the same method. First, identify the scientific ambiguity that the claim removes. Second, identify the typed OC construction that carries the distinction. Third, identify the proof sheet, formal fragment, finite semantic witness, replay row, comparator row, or claim-boundary record that can carry the public sentence. Fourth, identify the reopening condition. This common method is stated once here so that the following subsections can focus on the scientific work of each claim rather than repeat the same editorial apparatus.

The dependency route is layered rather than interchangeable. Definitions provide vocabulary; proof sheets state assumptions and lemmas; Lean records the encoded fragment; finite semantic cases show executable positive and negative behavior where the theorem is finite; empirical replay rows support only bounded artifact and reconstruction claims; comparator rows protect novelty wording; and review criteria check that the public sentence has not drifted beyond the evidence. No layer is allowed to impersonate another.

The hostile-review method is likewise fixed. A reviewer may remove a tuple component, collapse a status distinction, turn a boundary into a metaphor, treat a finite witness as a broad empirical prediction, or upgrade a comparator overlap into priority. A claim survives only while the named proof and evidence path blocks that attack at the stated strength. Otherwise the claim is repaired, demoted, or held for a later research release.

The editorial method is different from the proof method but it is not optional. The prose must let a reader see the problem, construction, evidence, and boundary without parsing a control table. Row identifiers serve as citations for audit and reproduction; they are not the voice of the scientific argument. The following subsections therefore preserve exact evidence identifiers while making each claim readable as science.

\subsection{Lifecycle status: liveness, death, residue, and rebirth}
The first question for this claim is what weaker vocabulary would confuse. The claim addresses this ambiguity: Death blocks live status; residue and rebirth are token-bound evidence relations with distinct class-specific endpoint rules: residue separates source from residue while returning to the source endpoint, and rebirth separates source, residue, and new target tokens. Rebirth is non-identity unless endpoint-bound identity evidence has identity class, declared invariant preservation, no residue token, and equal source/target endpoint evidence. No categorical Hom/composition theorem is promoted. The point of including it in the 1.3.3 release is to make the ambiguity auditable rather than to decorate the theory with another label.

The supporting material should be read through Lean-oriented formalization inventory, structured proof sheet, and finite witness. The proof route is public proof record for the lifecycle status theorem; the formal route is formalization reference lifecycle\_residue\_rebirth\_morphism\_boundary. In the public reading, the claim is carried by a named theorem boundary, not by the fact that a row exists. If the proof sheet, Lean reference, finite witness, or mutation control does not carry the visible wording, the visible wording is the part that must move.

The safe reading boundary is Bounded to stated theorem assumptions, finite witnesses, and public falsifier boundary. This boundary is not decorative. It prevents a bounded model-core statement from becoming a claim about unrestricted totality, full-domain prediction, or global priority. The scientific value of the claim is that it is strong enough to test a real structural distinction and narrow enough that an adversarial reviewer can name the condition that would reopen it.

A concrete reading starts by locating the typed field or operator that does the work, then looking for the positive case and the negative case that would have collapsed the distinction.

For T133-OMEGA-STATUS, the corresponding falsifier is concrete rather than slogan-like: find a model satisfying the stated assumptions in which the asserted distinction collapses, show that the positive witness was accepted for the wrong reason, show that the negative control is non-responsive, or show that the public wording requires an assumption absent from the proof sheet.

The prior-art and empirical boundaries remain local to this claim. Earlier formalisms may already carry part of the distinction; when they do, OC claims integration only for that part. Numeric replay rows can demonstrate artifact discipline, comparator behavior, and falsifier plumbing, but they do not by themselves complete a domain science. The reader should therefore leave with four handles: the conceptual problem, the typed construction, the evidence route, and the reopening route.

\subsection{Continuumness-zero obstruction}
The claim is useful only if it separates two situations that a less typed account would merge. The claim addresses this ambiguity: Continuumness zero requires live support, an independently clear obstruction register, and a declared zero-cause family; a zero-cause label alone does not compute k=0. The point of including it in the 1.3.3 release is to make the ambiguity auditable rather than to decorate the theory with another label.

The evidence ceiling for this claim is Lean-oriented formalization inventory, structured proof sheet, and finite witness. The proof route is public proof record for the K-zero boundary theorem; the formal route is formalization reference continuumness\_zero\_case\_iff\_declared\_zero\_cause\_with\_support. In the public reading, the claim is carried by a named theorem boundary, not by the fact that a row exists. If the proof sheet, Lean reference, finite witness, or mutation control does not carry the visible wording, the visible wording is the part that must move.

The claim stops at Bounded to stated theorem assumptions, finite witnesses, and public falsifier boundary. This boundary is not decorative. It prevents a bounded model-core statement from becoming a claim about unrestricted totality, full-domain prediction, or global priority. The scientific value of the claim is that it is strong enough to test a real structural distinction and narrow enough that an adversarial reviewer can name the condition that would reopen it.

The practical test is to ask what changes if the distinction is removed: which witness disappears, which counterexample becomes possible, and which comparator absorbs the claim.

For T133-K-ZERO, the corresponding falsifier is concrete rather than slogan-like: find a model satisfying the stated assumptions in which the asserted distinction collapses, show that the positive witness was accepted for the wrong reason, show that the negative control is non-responsive, or show that the public wording requires an assumption absent from the proof sheet.

The prior-art and empirical boundaries remain local to this claim. Earlier formalisms may already carry part of the distinction; when they do, OC claims integration only for that part. Numeric replay rows can demonstrate artifact discipline, comparator behavior, and falsifier plumbing, but they do not by themselves complete a domain science. The reader should therefore leave with four handles: the conceptual problem, the typed construction, the evidence route, and the reopening route.

\subsection{Metric-threshold boundary specialization}
The scientific burden here is to show why the named distinction is not merely verbal. The claim addresses this ambiguity: Metric thresholds are a specialization of typed classifier boundaries. The point of including it in the 1.3.3 release is to make the ambiguity auditable rather than to decorate the theory with another label.

The public warrant currently available is Lean-oriented formalization inventory, structured proof sheet, and finite witness. The proof route is public proof record for the boundary representation theorem; the formal route is formalization reference metric\_boundary\_specialization. In the public reading, the claim is carried by a named theorem boundary, not by the fact that a row exists. If the proof sheet, Lean reference, finite witness, or mutation control does not carry the visible wording, the visible wording is the part that must move.

The current release keeps the wording inside Bounded to stated theorem assumptions, finite witnesses, and public falsifier boundary. This boundary is not decorative. It prevents a bounded model-core statement from becoming a claim about unrestricted totality, full-domain prediction, or global priority. The scientific value of the claim is that it is strong enough to test a real structural distinction and narrow enough that an adversarial reviewer can name the condition that would reopen it.

The reader can inspect the claim by following the construction from named object to witness to reopening condition instead of treating the identifier as a proof.

For T133-BOUNDARY, the corresponding falsifier is concrete rather than slogan-like: find a model satisfying the stated assumptions in which the asserted distinction collapses, show that the positive witness was accepted for the wrong reason, show that the negative control is non-responsive, or show that the public wording requires an assumption absent from the proof sheet.

The prior-art and empirical boundaries remain local to this claim. Earlier formalisms may already carry part of the distinction; when they do, OC claims integration only for that part. Numeric replay rows can demonstrate artifact discipline, comparator behavior, and falsifier plumbing, but they do not by themselves complete a domain science. The reader should therefore leave with four handles: the conceptual problem, the typed construction, the evidence route, and the reopening route.

\subsection{Typed hybrid update and chart-labelled operator semantics}
This claim earns its public place by making a repair condition visible. The claim addresses this ambiguity: OC operators are typed update semantics; chart-labelled flow-one notation is admitted only for declared chart records, while proof/rewrite and guard/reset updates remain first-class non-smooth cases. No differentiability or ODE-solution theorem is promoted in declared formal profile. The point of including it in the 1.3.3 release is to make the ambiguity auditable rather than to decorate the theory with another label.

The reviewable support route is Lean-oriented formalization inventory, structured proof sheet, and finite witness. The proof route is public proof record for the hybrid semantics theorem; the formal route is formalization reference integrated\_operator\_semantics. In the public reading, the claim is carried by a named theorem boundary, not by the fact that a row exists. If the proof sheet, Lean reference, finite witness, or mutation control does not carry the visible wording, the visible wording is the part that must move.

The adversarial boundary is Bounded to stated theorem assumptions, finite witnesses, and public falsifier boundary. This boundary is not decorative. It prevents a bounded model-core statement from becoming a claim about unrestricted totality, full-domain prediction, or global priority. The scientific value of the claim is that it is strong enough to test a real structural distinction and narrow enough that an adversarial reviewer can name the condition that would reopen it.

A useful check is to translate the claim into one positive scenario and one failure scenario; if the failure scenario cannot be stated, the public wording is too strong.

For T133-HYBRID, the corresponding falsifier is concrete rather than slogan-like: find a model satisfying the stated assumptions in which the asserted distinction collapses, show that the positive witness was accepted for the wrong reason, show that the negative control is non-responsive, or show that the public wording requires an assumption absent from the proof sheet.

The prior-art and empirical boundaries remain local to this claim. Earlier formalisms may already carry part of the distinction; when they do, OC claims integration only for that part. Numeric replay rows can demonstrate artifact discipline, comparator behavior, and falsifier plumbing, but they do not by themselves complete a domain science. The reader should therefore leave with four handles: the conceptual problem, the typed construction, the evidence route, and the reopening route.

\subsection{Live-status cycle-mode requirement}
The reviewer-facing issue is whether the construction changes the space of admissible explanations. The claim addresses this ambiguity: Declared eligible-live status requires an explicit cycle mode or non-vacuous maintenance predicate. The point of including it in the 1.3.3 release is to make the ambiguity auditable rather than to decorate the theory with another label.

The claim is carried, if it is carried at all, by Lean-oriented formalization inventory, structured proof sheet, and finite witness. The proof route is public proof record for the cycle-mode theorem; the formal route is formalization reference eligible\_live\_requires\_cycle\_or\_maintenance. In the public reading, the claim is carried by a named theorem boundary, not by the fact that a row exists. If the proof sheet, Lean reference, finite witness, or mutation control does not carry the visible wording, the visible wording is the part that must move.

The public sentence is allowed only within Bounded to stated theorem assumptions, finite witnesses, and public falsifier boundary. This boundary is not decorative. It prevents a bounded model-core statement from becoming a claim about unrestricted totality, full-domain prediction, or global priority. The scientific value of the claim is that it is strong enough to test a real structural distinction and narrow enough that an adversarial reviewer can name the condition that would reopen it.

The local example is deliberately modest: it asks which boundary, morphism, status predicate, or operator prevents the claim from becoming a slogan.

For T133-CYCLE, the corresponding falsifier is concrete rather than slogan-like: find a model satisfying the stated assumptions in which the asserted distinction collapses, show that the positive witness was accepted for the wrong reason, show that the negative control is non-responsive, or show that the public wording requires an assumption absent from the proof sheet.

The prior-art and empirical boundaries remain local to this claim. Earlier formalisms may already carry part of the distinction; when they do, OC claims integration only for that part. Numeric replay rows can demonstrate artifact discipline, comparator behavior, and falsifier plumbing, but they do not by themselves complete a domain science. The reader should therefore leave with four handles: the conceptual problem, the typed construction, the evidence route, and the reopening route.

\subsection{Identity, residue, and rebirth classification}
The practical reading starts with the error that the claim prevents. The claim addresses this ambiguity: Identity continuation requires endpoint-bound identity evidence: identity class, declared invariant preservation, no residue token, equal source/target endpoint evidence, lifecycle identity-invariant truth, and typed source/target binding. Residue and rebirth evidence classes do not become identity continuation merely by preserving some invariants. No categorical Hom/composition theorem is promoted. The point of including it in the 1.3.3 release is to make the ambiguity auditable rather than to decorate the theory with another label.

The strongest visible support for this wording is Lean-oriented formalization inventory, structured proof sheet, and finite witness. The proof route is public proof record for the identity and rebirth theorem; the formal route is formalization reference endpoint\_bound\_identity\_classification. In the public reading, the claim is carried by a named theorem boundary, not by the fact that a row exists. If the proof sheet, Lean reference, finite witness, or mutation control does not carry the visible wording, the visible wording is the part that must move.

The reviewer should not read beyond Bounded to stated theorem assumptions, finite witnesses, and public falsifier boundary. This boundary is not decorative. It prevents a bounded model-core statement from becoming a claim about unrestricted totality, full-domain prediction, or global priority. The scientific value of the claim is that it is strong enough to test a real structural distinction and narrow enough that an adversarial reviewer can name the condition that would reopen it.

The subsection should leave the reader with an inspectable contrast between the case where the claim works and the case where the claim must be demoted.

For T133-ID, the corresponding falsifier is concrete rather than slogan-like: find a model satisfying the stated assumptions in which the asserted distinction collapses, show that the positive witness was accepted for the wrong reason, show that the negative control is non-responsive, or show that the public wording requires an assumption absent from the proof sheet.

The prior-art and empirical boundaries remain local to this claim. Earlier formalisms may already carry part of the distinction; when they do, OC claims integration only for that part. Numeric replay rows can demonstrate artifact discipline, comparator behavior, and falsifier plumbing, but they do not by themselves complete a domain science. The reader should therefore leave with four handles: the conceptual problem, the typed construction, the evidence route, and the reopening route.

\subsection{Declared component independence of the semantic verdict interface}
The first question for this claim is what weaker vocabulary would confuse. The claim addresses this ambiguity: Within the declared current release tuple semantics, each tuple component has a one-field semantic keep/drop witness that changes the release verdict. The point of including it in the 1.3.3 release is to make the ambiguity auditable rather than to decorate the theory with another label.

The supporting material should be read through Lean-oriented formalization inventory, structured proof sheet, and finite witness. The proof route is public proof record for the minimality witness theorem; the formal route is formalization reference release\_tuple\_semantic\_component\_irredundant. In the public reading, the claim is carried by a named theorem boundary, not by the fact that a row exists. If the proof sheet, Lean reference, finite witness, or mutation control does not carry the visible wording, the visible wording is the part that must move.

The safe reading boundary is Bounded to stated theorem assumptions, finite witnesses, and public falsifier boundary. This boundary is not decorative. It prevents a bounded model-core statement from becoming a claim about unrestricted totality, full-domain prediction, or global priority. The scientific value of the claim is that it is strong enough to test a real structural distinction and narrow enough that an adversarial reviewer can name the condition that would reopen it.

A concrete reading starts by locating the typed field or operator that does the work, then looking for the positive case and the negative case that would have collapsed the distinction.

For T133-MIN, the corresponding falsifier is concrete rather than slogan-like: find a model satisfying the stated assumptions in which the asserted distinction collapses, show that the positive witness was accepted for the wrong reason, show that the negative control is non-responsive, or show that the public wording requires an assumption absent from the proof sheet.

The prior-art and empirical boundaries remain local to this claim. Earlier formalisms may already carry part of the distinction; when they do, OC claims integration only for that part. Numeric replay rows can demonstrate artifact discipline, comparator behavior, and falsifier plumbing, but they do not by themselves complete a domain science. The reader should therefore leave with four handles: the conceptual problem, the typed construction, the evidence route, and the reopening route.

\subsection{Adjacent K-level witness consistency}
The claim is useful only if it separates two situations that a less typed account would merge. The claim addresses this ambiguity: Every declared adjacent K-level transition K0-\textgreater{}K12 has a release-atlas row, retained-witness evaluator check, executable finite row, and inert-witness demotion control inside the declared formal profile release classifier; independent semantic irreducibility beyond this declared classifier is a future proof obligation, not a promoted declared formal profile theorem. The point of including it in the 1.3.3 release is to make the ambiguity auditable rather than to decorate the theory with another label.

The evidence ceiling for this claim is Lean-oriented formalization inventory, structured proof sheet, and finite witness. The proof route is public proof record for the K-level witness theorem; the formal route is formalization reference release\_atlas\_manifest\_has\_total\_finite\_case\_coverage. In the public reading, the claim is carried by a named theorem boundary, not by the fact that a row exists. If the proof sheet, Lean reference, finite witness, or mutation control does not carry the visible wording, the visible wording is the part that must move.

The claim stops at Bounded to stated theorem assumptions, finite witnesses, and public falsifier boundary. This boundary is not decorative. It prevents a bounded model-core statement from becoming a claim about unrestricted totality, full-domain prediction, or global priority. The scientific value of the claim is that it is strong enough to test a real structural distinction and narrow enough that an adversarial reviewer can name the condition that would reopen it.

The practical test is to ask what changes if the distinction is removed: which witness disappears, which counterexample becomes possible, and which comparator absorbs the claim.

For T133-KLEVEL, the corresponding falsifier is concrete rather than slogan-like: find a model satisfying the stated assumptions in which the asserted distinction collapses, show that the positive witness was accepted for the wrong reason, show that the negative control is non-responsive, or show that the public wording requires an assumption absent from the proof sheet.

The prior-art and empirical boundaries remain local to this claim. Earlier formalisms may already carry part of the distinction; when they do, OC claims integration only for that part. Numeric replay rows can demonstrate artifact discipline, comparator behavior, and falsifier plumbing, but they do not by themselves complete a domain science. The reader should therefore leave with four handles: the conceptual problem, the typed construction, the evidence route, and the reopening route.


\section{How to Inspect the Integrated 1.3.3 Argument}
\label{sec:oc133-scientific-reading-protocol}
The integrated monograph is intentionally long because it has to remain useful to several kinds of external reader at once. The most reliable inspection path begins with the object grammar, moves to theorem and finite semantics, then to empirical replay, then to comparator boundaries, and finally to publication and citation surfaces. This order preserves the argument: a replay row is easier to judge after the reader knows the object whose behavior the row is meant to reconstruct.

The object grammar names carriers, realizations, lawful possibility, liveness, residue, identity morphisms, boundaries, operators, cycles, and K-level witnesses. The theorem layer then places a ceiling on wording: a promoted formal statement needs assumptions, proof route, finite witness interpretation, and a counterexample boundary. The executable layer contributes Lean-subset evidence and finite semantic cases where those checks exist. The empirical layer contributes retrospective replay rows and replay QA with formulas, inputs, comparators, residuals, negative controls, falsifiers, and hashes.

The comparator layer is equally important. OC does not benefit from pretending that systems theory, cybernetics, autopoiesis, dynamical systems, category theory, RAF theory, complexity measures, identity theory, or reproducibility engineering are absent. The useful question is narrower and more serious: which overlap is accepted, which residual delta remains, and which public statement is still justified under declared assumptions?

Reviewer objections are treated as scientific material when they name an attacked claim, identify the criticism route, point to the evidence that would answer it, and state the condition that would reopen the claim. This keeps criticism close to the model rather than turning it into a separate commentary layer.

The same principle governs repair. When a defect appears, the correction belongs at the first layer where the defect originates: ontology, theorem boundary, executable witness, empirical replay, comparator positioning, adversarial response, editorial prose, or publication packaging. Lower-layer errors should be fixed before a global reader is asked to judge the whole manuscript.
